% 第 2 讲 希尔伯特符号
% 源：src/lecture02.txt（Raskin, 18.786 讲义，2017-08-21 版）
% 说明：原文页眉「2. HILBERT SYMBOLS 7/9」为 PDF 页眉残留，已删。

\section*{第 2 讲\quad 希尔伯特符号}

设 $K$ 是 $\mathbb{Q}_p$ 上的局部域（其实任意局部域都够用），且 $\operatorname{char}(K)\neq 2$。注意这包含 $\mathbb{Q}_2$ 上的域，因为我们关心的是域的特征，而不是剩余域的特征。回忆上一讲的配对
\begin{equation}\tag{2.1}
  \operatorname{Gal}_2(K):=\operatorname{Gal}^{\mathrm{ab}}(K)/\{g^2:g\in\operatorname{Gal}^{\mathrm{ab}}(K)\}
  \cong\operatorname{Hom}\big(K^{\times}/(K^{\times})^2,\mathbb{Z}/2\mathbb{Z}\big),
\end{equation}
其中 $\operatorname{Gal}_2(K)$ 与 $\operatorname{Gal}^{\mathrm{ab}}(K)$ 是投射有限群，后者是 $K$ 的极大阿贝尔扩张的 Galois 群；而 $K^{\times}/(K^{\times})^2$ 是二元域 $\mathbb{Z}/2\mathbb{Z}$ 上有限维或无限维的向量空间（取对偶时，构成 $K^{\times}/(K^{\times})^2$ 的 $\mathbb{Z}/2\mathbb{Z}$ 直和变成了乘积，这反映了左端的投射有限性）。

还要回忆，lcft 断言 $K^{\times}\to\operatorname{Gal}^{\mathrm{ab}}(K)$ 是投射有限完备化，因此与 (2.1) 相反地有 $\operatorname{Gal}_2(K)\cong K^{\times}/(K^{\times})^2$。于是 lcft 预言存在如下形式的标准配对：

\textbf{定义 2.1.} 令希尔伯特符号
\[
  (\cdot,\cdot)\colon K^{\times}/(K^{\times})^2\times K^{\times}/(K^{\times})^2\to\{1,-1\}
\]
由
\[
  (a,b):=
  \begin{cases}
    1, & \text{若存在 } x,y\in K \text{ 使 } ax^2+by^2=1,\\
    -1, & \text{否则},
  \end{cases}
\]
定义，其中 $a,b\in K^{\times}$。

\textbf{注 2.2.} 这个定义只在局部域上性质良好。还要注意，$(a,b)$ 确实只在 $a,b$ 乘以平方的意义下定义，因为这些平方因子可以吸收进 $x$ 与 $y$。

\textbf{命题 2.3.} 希尔伯特符号满足以下性质：

(1) \emph{双乘性.} 对一切 $a,b,c\in K^{\times}$，
\[
  (a,bc)=(a,b)\cdot(a,c).
\]

(2) \emph{非退化性.} 对一切 $a\in K^{\times}$，若对一切 $b\in K^{\times}$ 都有 $(a,b)=1$，则 $a\in(K^{\times})^2$。

注意 $(a,b)=(b,a)$ 是平凡的。双乘性说的是：$ax^2+bcy^2=1$ 可解，当且仅当 $ax^2+by^2=1$ 与 $ax^2+cy^2=1$ 两者分别可解，或者两者都不可解。这有点奇怪，而且事实证明一般只在局部域上成立。
% [原文疑误：抽出文本两处均作 "ax²+by²=1"，按 (a,bc) 的含义第一个方程应为 ax²+bcy²=1；已按数学含义改正。]

\textbf{例 2.4.} 取 $K:=\mathbb{R}$。只要 $a$ 与 $b$ 不同时为负，$ax^2+by^2=1$ 就可解。因此 $\mathbb{R}^{\times}/(\mathbb{R}^{\times})^2=\{1,-1\}$（因为 $(\mathbb{R}^{\times})^2=\mathbb{R}_{>0}$），于是配对 $\{1,-1\}\times\{1,-1\}\to\{1,-1\}$ 确实非退化。

我们现在问：$x\in K^{\times}$ 何时是平方？当 $K=\mathbb{R},\mathbb{C}$ 时答案清楚。例如当 $x\in\mathbb{Q}_2^{\times}$ 时，可写 $x=2^{v(x)}y$，其中 $y\in\mathbb{Z}_2^{\times}$，而 $x$ 是平方当且仅当 $v_2(x)$ 为偶数且 $y$ 是平方（后者由习题集 1 第 1(c) 题将证明，等价于 $y\equiv 1\bmod 8$）。

设 $\mathfrak{p}\subseteq\mathcal{O}_K$ 为唯一的极大理想，$k:=\mathcal{O}_K/\mathfrak{p}$ 为剩余域，$\operatorname{char}(k)=p$ 为奇素数，$\pi\in\mathfrak{p}$ 为素元，即 $\pi\notin\mathfrak{p}^2$。

\textbf{断言 2.5.} 设 $x\in K^{\times}$，并写 $x=\pi^{v(x)}y$，其中 $y\in\mathcal{O}_K^{\times}$。则以下诸条等价：

(1) $x$ 是平方；

(2) $v(x)$ 为偶数，且 $y$ 是平方；

(3) $y\bmod\mathfrak{p}$ 是 $k^{\times}$ 中的平方。
% [原文疑误：抽出文本作 "a square in K^×"，按上下文应为 k^×（剩余域的单位群），见报告。]

注意我们可以化归到 $x\in\mathcal{O}_K^{\times}$。我们给出两个证明：

\textit{证明（用 Hensel 引理）.} 除最后一个条件外，其余解释都清楚。于是设 $x\bmod\mathfrak{p}$ 是 $\mathcal{O}_K^{\times}$ 中的平方。由 Hensel 引理，多项式 $p(t)=t^2-x\in\mathcal{O}_K[t]$ 有根 $r\in\mathcal{O}_K$，只要它有模 $\mathfrak{p}$ 的根 $\bar r$ 使 $\bar p'(\bar r)\neq 0$。第一个条件由假设成立：$x\bmod\mathfrak{p}=\bar r^2$ 对某个 $\bar r$ 成立。此时 $p'(t)=2t$，而 $\bar r\neq 0$（因 $\bar r^2=\bar x\neq 0$），又 $\operatorname{char}(k)\neq 2$，故 $\bar p'(\bar r)=2\bar r\neq 0$，于是第二个条件也成立。
\hfill$\square$

\textit{证明（显式）.} 考虑映射 $x\mapsto x^2$，即
\[
  \mathcal{O}_K^{\times}\xrightarrow{\ \sigma\ }S\subseteq\mathcal{O}_K^{\times},
  \qquad S:=\{x\in\mathcal{O}_K^{\times}:x \text{ 是模 } \mathfrak{p} \text{ 的平方}\}.
\]
我们想证明 $\sigma$ 满射，即 $\mathcal{O}_K$ 中每个模 $\mathfrak{p}$ 为平方的元素，都是 $\mathcal{O}_K^{\times}$ 中的平方。注意 $\mathcal{O}_K^{\times}$ 是带完全滤过的滤过阿贝尔群（见下文定义 2.7）
\[
  \mathcal{O}_K^{\times}\supseteq 1+\mathfrak{p}\supseteq 1+\mathfrak{p}^2\supseteq 1+\mathfrak{p}^3\supseteq\cdots,
\]
其中 $1+\mathfrak{p}^n$ 都是 $\mathcal{O}_K^{\times}$ 的开子群。显然 $\mathcal{O}_K^{\times}/(1+\mathfrak{p})=k^{\times}$；类似地 $(1+\mathfrak{p})/(1+\mathfrak{p}^2)\cong k$：对任意 $1+a\pi,1+b\pi\in 1+\mathfrak{p}$（$a,b\in\mathcal{O}_K/\mathfrak{p}$），有
\[
  (1+a\pi)(1+b\pi)=1+(a+b)\pi+ab\pi^2,
\]
而 $ab\pi^2\in\mathfrak{p}^2$，故在相伴分次项中剩下 $1+(a+b)\pi$，于是乘法恰好对应 $k$ 中的加法。类似地，对每个 $n\geq 1$，由同样的论证可得 $(1+\mathfrak{p}^n)/(1+\mathfrak{p}^{n+1})\cong k$，因为 $n+1\leq 2n$。现在 $\sigma$ 如下作用在滤过上：
\[
  \begin{CD}
    \mathcal{O}_K^{\times} @<<< 1+\mathfrak{p} @<<< 1+\mathfrak{p}^2 @<<< \cdots\\
    @V{\sigma}VV @V{\sigma}VV @V{\sigma}VV\\
    \mathcal{O}_K^{\times} @<<< S @<<< 1+\mathfrak{p} @<<< \cdots
  \end{CD}
\]
其中包含 $1+\mathfrak{p}\subseteq S$ 成立，因为 $1$ 平凡地是平方。现在，$\mathrm{Gr}_0$ 上的映射
\[
  \mathcal{O}_K^{\times}/(1+\mathfrak{p})=k^{\times}\xrightarrow{\ \sigma\ }(k^{\times})^2=S/(1+\mathfrak{p})
\]
平凡地满射（核很小）。此外，对每个 $n\geq 1$，$\mathrm{Gr}_n$ 上的映射是
\[
  (1+\mathfrak{p}^n)/(1+\mathfrak{p}^{n+1})\xrightarrow{\ \sigma\ }(1+\mathfrak{p}^n)/(1+\mathfrak{p}^{n+1}),
\]
由于对任意 $x\in k$ 有
\[
  (1+\pi^n x)^2=1+2x\pi^n+x^2\pi^{2n}\equiv 1+2x\pi^n \pmod{\mathfrak{p}^{n+1}}
\]
（因 $2n\geq n+1$），它等同于映射 $k\to k$，$x\mapsto 2x$；而 $\#k=p$ 为奇素数，故这是同构。于是 $\sigma$ 在每个分次项上都满射，由命题 2.9 便知 $\mathcal{O}_K^{\times}\xrightarrow{\sigma}S$ 满射，如所欲证。
\hfill$\square$

\textbf{注 2.6.} 一般说来，我们处理 $\mathcal{O}_K^{\times}$ 的工具是 $p$ 进指数映射与这个滤过；后者虽然只是抽象的形式体系，却比 $\mathcal{O}_K^{\times}$ 更简单，因为其商全都同构于有限域。作为一般原则，我们可以通过相伴分次 $\mathrm{Gr}_*A$ 来理解关于 $A$ 的许多事情。

\textbf{定义 2.7.} 设 $A$ 为阿贝尔群。$A$ 上的一个\emph{滤过}是一列递降的子群
\[
  A=:F^0A\supseteq F^1A\supseteq F^2A\supseteq\cdots,
\]
若 $A\xrightarrow{\sim}\varprojlim_n A/F^nA$，则称该滤过\emph{完全}。群 $\mathrm{Gr}_nA:=F^nA/F^{n+1}A$ 称为滤过的\emph{相伴分次项}。

\textbf{例 2.8.} 群
\[
  \mathcal{O}_K\cong\varprojlim_n \mathcal{O}_K/\mathfrak{p}^n,
  \qquad
  \mathcal{O}_K^{\times}\cong\varprojlim_n \mathcal{O}_K^{\times}/(1+\mathfrak{p}^n)
\]
都是完全滤过。

\textbf{命题 2.9.} 设 $f\colon A\to B$ 是完全滤过阿贝尔群之间的同态，即对每个 $n\geq 0$ 有 $f(F^nA)\subset F^nB$。若诱导映射
\[
  F^nA/F^{n+1}A\to F^nB/F^{n+1}B
\]
满射（相应地，单射），则 $f$ 满射（相应地，单射）。

\textit{证明.} 假设相伴分次上的映射都满射，且两个滤过都完全，如断言 2.5 的显式证明中那样。设有 $x\in B$，我们想解方程 $f(y)=x$，$y\in A$。先可解 $f(y_0)\equiv x\bmod F^1B$，于是 $x-f(y_0)\in F^1B$。其次，由假设相伴分次映射满射，可解 $f(\varepsilon_1)\equiv x-f(y_0)\bmod F^2B$，其中 $\varepsilon_1\in F^1A$ 是从 $\mathrm{Gr}_1A$ 提升来的「误差项」。注意由 $f$ 是同态，有
\[
  f(y_1)=f(y_0+\varepsilon_1)=f(y_0)+f(\varepsilon_1)\equiv x \pmod{F^2B},
\]
其中我们记 $y_1:=y_0+\varepsilon_1$。这与之前的方程形式相同，故可迭代得到「相容」的 $y_n$，使 $f(y_n)\equiv x\bmod F^{n+1}B$ 对每个 $n\geq 0$ 成立；这里「相容」指 $y_n\equiv y_{n+1}\bmod F^{n+1}A$。但于是有诱导元素 $y\in\varprojlim A/F^nA=A$ 对应 $(y_0,y_1,\ldots)$（注意 $y_n$ 在 $n$ 充分大时模 $F^nA$ 稳定），它满足原方程 $f(y)=x$，因为两个滤过都由假设完全。\hfill$\square$

\textbf{注 2.10.} 尽管简单而抽象，许多事情（如上一个断言）都可用上述命题轻松证明。用 Hensel 引理的办法的优点在于：这里我们需要用平方映射 $\sigma$ 是同态这一事实，而这一般不成立。尽管如此，这个办法还是能告诉我们，在剩余特征为奇的局部域中 $\mathcal{O}_K^{\times}$ 的哪些元素是平方。

结论是：当 $K$ 是剩余特征为奇的局部域时，有 $[K^{\times}:(K^{\times})^2]=4$，因为 $[\mathcal{O}_K^{\times}:(\mathcal{O}_K^{\times})^2]=2$，而 $2\mathbb{Z}\subseteq\mathbb{Z}$ 同理；于是 $K^{\times}/(K^{\times})^2$ 同构于 $\mathbb{Z}/2\mathbb{Z}\times\mathbb{Z}/2\mathbb{Z}$，它以 $\{\pi,r\}\subset\mathcal{O}_K^{\times}$ 为基，其中 $\pi$ 是素元，$r$ 不是模 $\mathfrak{p}$ 的平方（因此 $\pi$ 与 $r$ 当然不会相差一个平方）。

现在我们把希尔伯特符号用扩张上的范数重新表述；与最初的定义相反，这里我们将不对称地看待它。设 $a$ 不是平方，于是 $K(\sqrt{a})$ 是 $K$ 的二次扩张（注意若 $a$ 是非零平方，则只需把 $K(\sqrt{a})$ 理解为 $K$ 的相应平展扩张，同构于 $K\times K$）。

\textbf{断言 2.11.} $(a,b)=1$ 当且仅当 $b$ 是扩张 $K(\sqrt{a})/K$ 的范数，即 $K(\sqrt{a})$ 中有元素其范数为 $b$。

\textit{证明.} 设 $b$ 是范数，即存在 $\alpha,\beta\in K$ 使
\[
  \alpha^2-\beta^2a=N(\alpha+\beta\sqrt{a})=b,
\]
从而 $\alpha^2=a\beta^2+b$。若 $\alpha\neq 0$，则
\[
  a\Big(\frac{\beta}{\alpha}\Big)^2+b\Big(\frac{1}{\alpha}\Big)^2=1,
\]
故 $(a,b)=1$。若 $\alpha=0$，则 $b+\beta^2a=0$。令
\[
  x:=\frac{1}{2}\Big(1+\frac{1}{a}\Big),
  \qquad
  y:=\frac{1}{2\beta}\Big(1-\frac{1}{a}\Big),
\]
我们有
\[
  ax^2+by^2
  =a\cdot\frac{1}{4}\cdot\frac{(a+1)^2}{a^2}
  +(-\beta^2a)\cdot\frac{1}{4\beta^2}\cdot\frac{(a-1)^2}{a^2}
  =\frac{(a+1)^2-(a-1)^2}{4a}=1,
\]
故仍有 $(a,b)=1$。反向蕴涵只是 $\alpha\neq 0$ 时论证的平凡倒推。\hfill$\square$

我们不加证明地陈述关于希尔伯特符号的主要结果。重要的是：$L^{\times}$ 在范映射下的像不能太大（不是全部），也不能太小。我们将看到，这个定理等价于希尔伯特符号的非退化性。

\textbf{定理 2.12.} 若 $L/K$ 是局部域的二次扩张，则范映射 $N\colon L^{\times}\to K^{\times}$ 是同态，且 $N(L^{\times})\subseteq K^{\times}$ 是指数为 $2$ 的子群。

\textbf{例 2.13.} 考虑 $\mathbb{C}/\mathbb{R}$。
% [原文在 src 抽取中于此句后截断（原 PDF 该例的后续文字未进入文本层），此处照存。]
